\documentclass[11pt,a4paper]{article}
\usepackage[T1]{fontenc}
\usepackage[utf8]{inputenc}
\usepackage[french]{babel}
\usepackage{lmodern}
\usepackage{microtype}
\usepackage[a4paper,margin=2.2cm]{geometry}
\usepackage{xcolor}
\usepackage{booktabs,tabularx,array}
\usepackage{enumitem}
\usepackage{hyperref}
\usepackage{fancyhdr}
\usepackage{titlesec}
\usepackage{tcolorbox}
\hypersetup{colorlinks=true,linkcolor=blue!40!black,urlcolor=blue!40!black}
\definecolor{deepblue}{HTML}{183B56}
\definecolor{softblue}{HTML}{EAF2F8}
\definecolor{softamber}{HTML}{FFF4D6}
\definecolor{softred}{HTML}{FBE9E7}
\definecolor{softgreen}{HTML}{E8F5E9}
\definecolor{muted}{HTML}{5F6B76}
\pagestyle{fancy}
\fancyhf{}
\lhead{ED-POMDP / STRAT-Q}
\rhead{Note d'état Step 2.7}
\cfoot{\thepage}
\setlength{\parindent}{0pt}
\setlength{\parskip}{0.6em}
\setlist[itemize]{leftmargin=1.5em,itemsep=0.25em,topsep=0.25em}
\titleformat{\section}{\large\bfseries\color{deepblue}}{}{0pt}{}
\titleformat{\subsection}{\normalsize\bfseries\color{deepblue}}{}{0pt}{}
\newcommand{\claimvoi}{\texttt{CLM-VOI-001}}
\newcommand{\claimeq}{\texttt{CLM-EQ-001}}

\begin{document}

{\LARGE\bfseries\color{deepblue} ED-POMDP -- Note d'état exécutive après Step 2.7}\par
{\large Décider avec des preuves, sans confondre risque du système et faiblesse des preuves}\par
{\color{muted}Version du 31 juillet 2026 -- Résultats gelés, exécutés et vérifiés indépendamment}

\begin{tcolorbox}[colback=softblue,colframe=deepblue,boxrule=0.7pt,arc=2mm]
\textbf{En une phrase.} Le modèle ED-POMDP produit souvent des probabilités mieux calibrées, mais cette amélioration ne se transforme pas en meilleures décisions finales dans le benchmark actuel. Le pari central de supériorité décisionnelle n'est donc pas soutenu.
\end{tcolorbox}

\section*{Ce qu'on cherchait}

Lorsqu'une organisation décide de livrer une solution -- \textbf{GO}, \textbf{CONDITIONAL GO} ou \textbf{NO-GO} -- elle observe des tests, des couvertures, des incidents, des traces ou des indicateurs d'environnement. Les approches classiques confondent souvent deux situations pourtant très différentes :

\begin{itemize}
  \item le système est réellement risqué ou défaillant ;
  \item les preuves disponibles sont elles-mêmes faibles, obsolètes, incomplètes ou peu représentatives.
\end{itemize}

ED-POMDP cherche à représenter séparément l'état réel du système et la qualité des preuves. L'objectif est aussi de déterminer \textbf{quelle preuve acquérir ensuite}, au lieu d'accumuler des tests selon un plan fixe ou par opportunité. Cette question est directement liée au rôle de Test Authority et à la logique STRAT-Q : réduire l'incertitude décisionnelle avec un budget de preuve limité.

\section*{Ce qu'on espérait démontrer}

Deux paris ont été posés comme des hypothèses réfutables, et non comme des vérités :

\begin{itemize}
  \item \textbf{Pari 1 -- valeur de l'information :} choisir intelligemment la prochaine preuve devait améliorer la décision terminale à budget égal, par rapport à des stratégies fixes, aléatoires, entropiques ou uniquement orientées risque (\claimvoi).
  \item \textbf{Pari 2 -- qualité de la preuve :} modéliser explicitement la qualité des preuves devait améliorer la calibration des probabilités et, idéalement, la décision finale lorsque les preuves sont dégradées ou le modèle mal spécifié (\claimeq).
\end{itemize}

\section*{Ce qui a été évalué}

Sept politiques de décision ont été comparées sur \textbf{3 360 épisodes simulés}, dans quatre régimes :

\begin{itemize}
  \item preuve identifiable et globalement fiable ;
  \item preuve dégradée ;
  \item vraisemblance mal spécifiée ;
  \item cas non identifiable, où le système et la qualité des preuves ne peuvent pas être séparés de manière fiable.
\end{itemize}

Quatre budgets d'acquisition ont été utilisés, avec 30 seeds gelés et les mêmes scénarios latents pour toutes les politiques. Les règles, endpoints, paramètres, nombres de bootstrap et de permutations ainsi que la correction de multiplicité ont été verrouillés avant l'exécution. Le reviewer a ensuite vérifié les hashes, les comptes de lignes et recalculé indépendamment Holm step-down avec \textbf{zéro écart sur 240 valeurs ajustées}.

\section*{Le résultat}

Sur les \textbf{240 comparaisons confirmatoires} prévues, une seule résiste à la correction de Holm : ED-POMDP améliore l'erreur de calibration ECE face au POMDP classique dans le régime de preuve dégradée, au budget 2. Le résultat est étroit : l'intervalle bootstrap de la différence \textbf{traverse zéro} et la cellule ne contient que peu de valeurs de postérieur distinctes.

\begin{tcolorbox}[colback=softamber,colframe=orange!70!black,boxrule=0.7pt,arc=2mm]
\textbf{Point central.} Sur la décision finale, l'absence de résultat significatif ne signifie pas seulement « on ne sait pas ». Sur les 64 comparaisons directement visées par le pari VOI :
\begin{itemize}
  \item ED-POMDP est meilleur dans 10 cellules ;
  \item moins bon dans 16 cellules ;
  \item exactement égal dans 38 cellules.
\end{itemize}
Ainsi, ED-POMDP \textbf{n'améliore pas} la décision dans 54 cas sur 64. Ce total combine beaucoup d'égalités et une direction défavorable parmi les cellules où les méthodes diffèrent ; il ne faut donc pas le traduire par « ED est moins bon quatre fois sur cinq ».
\end{tcolorbox}

La dissociation entre qualité probabiliste et qualité décisionnelle est nette :

\begin{center}
\begin{tabularx}{\textwidth}{>{\bfseries}l r r r X}
\toprule
Endpoint & Favorable & Défavorable & Égal & Lecture descriptive \\
\midrule
Perte de décision & 17/80 & 17/80 & 46/80 & Beaucoup d'égalités ; sur les quatre baselines VOI, les directions défavorables dominent 16 à 10. \\
Brier score & 63/80 & 13/80 & 4/80 & Les probabilités d'ED-POMDP sont plus souvent proches de la vérité. \\
ECE & 53/80 & 23/80 & 4/80 & La calibration est souvent meilleure, mais une seule cellule est confirmée après Holm. \\
\bottomrule
\end{tabularx}
\end{center}

En clair, ED-POMDP semble parfois \textbf{mieux estimer} le risque, mais cette meilleure estimation ne modifie pas assez souvent l'action terminale, ou ne la modifie pas dans la bonne direction. Il « pense » souvent mieux sans décider globalement mieux.

\section*{Ce que cela invalide ou limite}

\begin{tcolorbox}[colback=softred,colframe=red!60!black,boxrule=0.7pt,arc=2mm]
\textbf{Pari 1 -- non soutenu.} Le benchmark ne démontre pas qu'une acquisition décisionnelle de preuve produit de meilleures décisions finales à budget égal. La direction descriptive sur les baselines directement visées par le claim est même préoccupante : parmi les cellules non égales, les résultats défavorables sont plus nombreux que les favorables.
\end{tcolorbox}

\begin{tcolorbox}[colback=softgreen,colframe=green!45!black,boxrule=0.7pt,arc=2mm]
\textbf{Pari 2 -- soutien très étroit uniquement.} Une amélioration confirmatoire de calibration apparaît dans une seule configuration de preuve dégradée et de petit budget. Cela ne soutient ni une supériorité générale sous mauvaise spécification, ni une amélioration générale de la décision terminale.
\end{tcolorbox}

Les claims larges \claimvoi{} et \claimeq{} restent donc non promus. Le résultat négatif est une contribution scientifique réelle : il empêche de transformer une intuition séduisante en promesse non démontrée.

\section*{La suite : Step 2.8}

Step 2.8 doit formaliser le résultat et expliquer le mécanisme. L'hypothèse principale à examiner est que la calibration et la décision terminale sont insuffisamment couplées : améliorer la croyance ne sert pas si la règle finale, les seuils, l'horizon ou la fonction de perte n'exploitent pas cette amélioration.

Trois chemins sont désormais ouverts :

\begin{itemize}
  \item \textbf{revoir le mécanisme de décision} : règle terminale, seuils, politique VOI, coûts asymétriques et manière dont la qualité de preuve influence l'action ;
  \item \textbf{tester sur des cas plus réalistes} : flux STRAT-Q, preuves corrélées, environnements représentatifs, coûts et décisions GO / CONDITIONAL GO / NO-GO ;
  \item \textbf{documenter proprement la limite} : réfuter ou restreindre le claim VOI pour ce benchmark et préserver les résultats nuls, défavorables et de sécurité dans le papier.
\end{itemize}

\begin{tcolorbox}[colback=softblue,colframe=deepblue,boxrule=0.7pt,arc=2mm]
\textbf{État de décision scientifique.} Step 2.7 ne valide pas la supériorité générale d'ED-POMDP. Il révèle plutôt une question de recherche plus précise et plus utile : \emph{dans quelles conditions une meilleure représentation de la qualité des preuves change-t-elle réellement une décision d'ingénierie, et pas seulement la calibration d'une probabilité ?}
\end{tcolorbox}

\vfill
{\small\color{muted}Sources internes : protocole gelé Step 2.6, résultats immuables Step 2.7, audit SHA-256, analyse confirmatoire de 240 contrastes et revue indépendante. Cette note est une synthèse exécutive ; les tables détaillées restent dans le paquet scientifique.}

\end{document}
